\subsection{Системы координат и единицы}
\label{sec:sdk2-frames}

Команды точек работают в двух системах координат. \emph{Локальная} система неподвижна: начало $(0,0)$ --- точка старта, ось $Z$ направлена вверх. \emph{Связанная} (body fixed) система движется вместе с дроном: $Y_b$ --- вперёд (куда смотрит камера), $X_b$ --- вправо.

\begin{center}
\begin{tikzpicture}[>=latex, scale=0.9]
  \draw[->, thick] (0,0) -- (5.6,0) node[below left] {$X$};
  \draw[->, thick] (0,0) -- (0,4.0) node[below left] {$Y$};
  \fill (0,0) circle (1.5pt) node[below right, font=\small\sffamily] {старт $(0,0)$};
  \draw[densely dotted, gray] (3.2,0) node[below, font=\small] {$x$} -- (3.2,2.2) -- (0,2.2) node[left, font=\small] {$y$};
  \begin{scope}[shift={(3.2,2.2)}, rotate=30]
    \foreach \a in {45,135,225,315} \draw[GeoskanBlue, thick] (0,0) -- (\a:0.55) (\a:0.55) circle (0.18);
    \draw[fill=GeoskanBlue!15, draw=GeoskanBlue, thick] (-0.22,-0.22) rectangle (0.22,0.22);
    \draw[->, very thick, red!70!black] (0,0) -- (0,1.5) node[above, font=\small] {$Y_b$ вперёд};
    \draw[->, very thick, teal] (0,0) -- (1.5,0) node[right, font=\small] {$X_b$ вправо};
  \end{scope}
\end{tikzpicture}
\end{center}

\verb|go_to_local_point(x, y, z, yaw)| --- абсолютная точка в неподвижной системе; \verb|go_to_local_point_body_fixed(...)| --- смещение в осях корпуса от текущего положения. Единицы и знаки:

\begin{center}
\small
\begin{tabularx}{\linewidth}{@{}>{\raggedright\arraybackslash}p{0.34\linewidth} X@{}}
\toprule
\textbf{Величина} & \textbf{Единицы и смысл} \\
\midrule
\verb|x, y, z| & метры; $Z$ вверх \\
\verb|yaw| в \verb|go_to_*| & \textbf{градусы} (см. предупреждение ниже) \\
\verb|set_yaw(yaw)| & градусы; по описанию Pioneer Code положительный угол --- по часовой стрелке \\
\verb|vx, vy, vz| & м/с; в \verb|*_body_fixed|: $v_y > 0$ --- вперёд, $v_x > 0$ --- вправо \\
\verb|yaw_rate| & рад/с; положительное значение --- поворот влево (против часовой) \\
\verb|led_control(r, g, b)| & яркость каналов от 0 до 1; \verb|led_id=255| --- все светодиоды \\
\bottomrule
\end{tabularx}
\end{center}

\begin{warnbox}[Единицы курса: проверьте на своём дроне]
В исходниках 0.11.0, в справочнике симулятора и в программах этой главы курс \verb|yaw| в \verb|go_to_local_point| и \verb|go_to_local_point_body_fixed| --- в градусах. В сборнике примеров Pioneer Code встречается и другое утверждение: что в этих командах курс в радианах (примеры считают его через \verb|math.atan2| без перевода). Расхождение не удалось снять без полёта, поэтому проверьте единицы сами: дважды вызовите \verb|go_to_local_point(0, 0, 1, yaw=90)| и \verb|yaw=1.57|, после каждого печатайте \verb|get_local_yaw_lps()|. Если первый вызов даёт курс около 90, а второй --- около 1.6, то курс в градусах; если первый поворачивает дрон на странный угол (90 рад $\approx 117^\circ$ после приведения к диапазону $\pm180$), а второй даёт около 90 --- в радианах. Знак поворота тоже сверьте глазами.
\end{warnbox}

\begin{infobox}[Навигация]
Локальные координаты имеют смысл, только если работает навигация: \texttt{OPT} (оптический поток --- типично для помещений), \texttt{LPS} (локальная система позиционирования) или \texttt{GPS}. Активную систему показывает \verb|get_nav_system()|, а качество LPS и GPS --- \verb|get_nav_status_lps()| и \verb|get_nav_status_gps()| (\texttt{NO\_DATA}, \texttt{CANNOT}, \texttt{LOW}, \texttt{OK}); подробнее --- в разделе про навигацию.
\end{infobox}

\subsection{Полёт по точкам}

\begin{codefile}[python]{python-onboard/pioneer_sdk2/sdk2_03_square.py}
\end{codefile}

\paragraph{Что происходит.} \verb|go_to_local_point| в SDK~2 \emph{не} блокирующая: команда отправляет точку и сразу возвращается. Дожидаться прибытия по-прежнему нужно через \verb|point_reached()|, как в главе~\ref{sec:python-sdk}. Курс --- в \textbf{градусах}: 0 --- нос вдоль оси X, 90 --- вдоль оси Y. В списке у каждой вершины свой курс, и дрон летит по квадрату носом вперёд. Отдельной точки над стартом здесь нет: \verb|takeoff()| возвращается уже после взлёта.

\begin{codebox}[Команды точек]
\verb|go_to_local_point(x, y, z, yaw=None, time=0)| --- лететь в точку локальной системы координат (метры). \verb|yaw| --- курс в градусах; если не указан, курс не меняется. \verb|time| --- желаемое время перелёта в секундах; 0 --- лететь с текущей скоростью.\\
\verb|go_to_local_point_body_fixed(x, y, z, yaw=None, time=0)| --- то же, но координаты --- смещение относительно текущего положения дрона.\\
\verb|point_reached()| --- \verb|True| один раз после прибытия, затем снова \verb|False|.\\
\verb|point_deceleration()| --- \verb|True| один раз, когда дрон начал тормозить перед точкой.
\end{codebox}

Обе команды точек возвращают \verb|False|, если SDK проигнорировал их по таблице состояний (дрон ещё на земле); в нормальном полёте возвращаемого значения нет.

\begin{infobox}[Смещение «от дрона»: в какую сторону оси]
В симуляторе у \verb|go_to_local_point_body_fixed| ось $x$ направлена вправо от носа дрона, $y$ --- вперёд, $z$ --- вверх, а \verb|yaw| --- поворот от текущего курса в градусах. Поэтому «лететь вперёд на 1 м» --- это \verb|y=1|, а не \verb|x=1|. Как ориентированы оси смещения на вашем борту, проверьте коротким полётом, прежде чем строить на этом маршрут.
\end{infobox}

\begin{rulebox}[Градусы и радианы]
В SDK~2 курс \verb|yaw| в \verb|go_to_local_point| и \verb|set_yaw| --- в градусах, а скорость поворота \verb|yaw_rate| в \verb|set_manual_speed| --- в радианах в секунду. Переносите программы из главы~\ref{sec:python-sdk} внимательно: \verb|yaw=math.pi| там --- это \verb|yaw=180| здесь.
\end{rulebox}

\paragraph{Торможение перед точкой.} \verb|point_deceleration()| срабатывает раньше, чем \verb|point_reached()|: автопилот сообщает, что дрон подлетает и начал сбрасывать скорость. Это удобное место, чтобы заранее подготовить следующий шаг --- например, включить подсветку или навести камеру. Пример \texttt{sdk2\_09\_safe\_start.py} в конце раздела про навигацию использует оба признака.

\begin{warnbox}[Опрашивайте каждую точку]
SDK запоминает события, даже если вы их не опрашиваете, и \verb|point_reached()| сбрасывает флаг только когда вернул \verb|True|. Если вы отправили две точки подряд и ни разу не вызвали \verb|point_reached()| для первой, флаг «прибыл» может остаться от неё, и опрос второй точки сразу вернёт \verb|True|. Правило: после каждой точки --- свой цикл ожидания.
\end{warnbox}

\begin{taskbox}[Вперёд и назад в градусах]
Перенесите программу «Вперёд и назад» из главы~\ref{sec:python-sdk} на SDK~2: курс в градусах, ожидание точки через \verb|point_reached()| с тайм-аутом. На обратном пути летите носом к старту (\verb|yaw=180|).
\expected{дрон пролетает 1~м по $+X$, разворачивается на $180^\circ$, возвращается и садится на старте.}
\end{taskbox}

\begin{taskbox}[Время перелёта]
Отправьте дрона в точку $(2, 0, 1)$ с \verb|time=0|, замерьте время до \verb|point_reached()| часами \verb|time.time()|; потом то же самое с \verb|time=8|. Заодно запомните, через сколько секунд после команды срабатывает \verb|point_deceleration()|.
\expected{при \texttt{time=0} перелёт идёт с обычной скоростью и занимает заметно меньше 8~с; при \texttt{time=8} время перелёта близко к 8~с (точность зависит от дрона). Торможение срабатывает перед самым прибытием.}
\end{taskbox}

\subsubsection*{Миссия как данные: список точек}

\begin{codefile}[python]{python-onboard/pioneer_sdk2/05_square_mission.py}
\end{codefile}

\paragraph{Что происходит.} Маршрут --- это \emph{данные}: список \verb|WAYPOINTS| отделён от логики полёта, и чтобы облететь треугольник или «змейку», достаточно поменять список. \verb|fly_to()| отправляет неблокирующую команду и опрашивает \verb|point_reached()| каждые 0.1~с. \verb|time.monotonic()| --- часы, которые не скачут при смене системного времени, ими удобно мерить интервалы; если точка не достигнута за \verb|POINT_TIMEOUT|, функция возвращает \verb|False|, и миссия прекращается. Светодиоды показывают фазу: синий --- в пути, зелёный --- на точке.

\begin{infobox}[Аргумент time]
Последний аргумент \verb|go_to_local_point(..., time)| --- время, за которое нужно достичь точки. Значение 0 оставляет выбор скорости автопилоту; положительное значение делает перелёт предсказуемым по времени.
\end{infobox}

\begin{taskbox}[Змейка]
Сгенерируйте «змейку» для обследования прямоугольника $2\times1$~м с шагом 0.5~м функцией, а не вручную. Добавьте к точке четвёртый элемент --- курс --- и разверните дрон носом по направлению движения. Выведите длину маршрута и фактическое время полёта.
\expected{дрон проходит змейку строка за строкой; в журнале --- длина маршрута и время полёта.}
\end{taskbox}

\subsubsection*{Относительные перемещения}

\begin{codefile}[python]{python-onboard/pioneer_sdk2/06_body_fixed_moves.py}
\end{codefile}

\paragraph{Что происходит.} \verb|go_to_local_point_body_fixed()| задаёт \emph{смещение} относительно текущего положения и курса: «вперёд» --- всегда туда, куда смотрит камера. Функция \verb|step()| переименовывает оси в понятные слова: \verb|forward| $\to$ \verb|y|, \verb|right| $\to$ \verb|x|. Четыре пары «вперёд + поворот» замыкают квадрат, хотя программа не знает ни одной абсолютной координаты.

\begin{taskbox}[Шестиугольник и спираль]
Нарисуйте в воздухе правильный шестиугольник (угол поворота $360^\circ/6$). Напишите функцию \verb|spiral(n)|: каждый шаг вперёд на 0.2~м длиннее предыдущего. Сравните: какой из примеров, 05 или 06, точнее возвращается в точку старта? Почему ошибки в относительном режиме накапливаются?
\expected{шестиугольник замыкается; спираль раскручивается; относительный маршрут заканчивается дальше от старта, потому что ошибка каждого шага прибавляется к следующему.}
\end{taskbox}

\subsubsection*{Облёт по окружности: камера смотрит в центр}

\begin{codefile}[python]{python-onboard/pioneer_sdk2/07_circle_face_center.py}
\end{codefile}

\paragraph{Что происходит.} Точка на окружности радиуса $R$ с центром $(x_c, y_c)$ задаётся углом $\varphi_i = 2\pi i/N$:
\[
x = x_c + R\cos\varphi,\qquad y = y_c + R\sin\varphi.
\]
Чтобы нос смотрел в центр, курс должен совпадать с направлением вектора из точки в центр: $\psi = \operatorname{atan2}(y_c - y,\; x_c - x)$. Функция \verb|math.atan2(dy, dx)|, в отличие от \verb|atan(dy/dx)|, учитывает знаки и работает во всех четырёх четвертях. Цикл идёт до \verb|POINTS + 1|, чтобы последняя точка совпала с первой. Чем больше \verb|POINTS|, тем ближе многоугольник к окружности, но тем чаще дрон тормозит в точках. Курс в примере переведён в градусы (\verb|math.degrees|), как и во всей главе; см. предупреждение про единицы курса выше.

\begin{taskbox}[Спираль и восьмёрка]
Сделайте спираль: радиус уменьшается на 0.05~м на каждой точке, а высота растёт. Нарисуйте «восьмёрку» (лемнискату): $x = R\sin\varphi$, $y = R\sin\varphi\cos\varphi$.
\expected{спираль сходится к центру и поднимается; восьмёрка пересекает центр дважды за круг.}
\end{taskbox}

\subsubsection*{Глобальные точки по GPS}

Локальные координаты отсчитываются в системе локальной навигации (LPS). Если на дроне установлен модуль GPS/ГНСС, то летать можно по географическим координатам:

\begin{codebox}[Глобальные точки (по 0.11.0)]
\verb|go_to_global_point(latitude, longitude, altitude, yaw=0)| --- лететь в точку: широта и долгота в градусах, высота над уровнем моря в метрах, курс в градусах.\\
\verb|go_to_global_point_relative(latitude_offset, longitude_offset, altitude_offset, yaw=0)| --- лететь в точку, заданную смещением от текущего положения.
\end{codebox}

Обе команды работают только когда навигация дрона --- GPS: если \verb|get_nav_system()| вернул не \verb|GPS|, библиотека печатает «Команда может быть выполнена только когда установлена система навигации GPS» и возвращает \verb|False|. На земле команда игнорируется (\verb|False|), в состоянии \verb|ARMED| --- \verb|RuntimeError|, как у любой полётной команды. Результат \verb|True| означает «команда отправлена», а не «дрон прилетел»: прибытие определяется так же, как для локальных точек --- через \verb|point_reached()|.

\begin{codebox}[Пример: сместиться на север (по 0.11.0)]
\small
\begin{verbatim}
from pioneer_sdk2 import Pioneer

pioneer = Pioneer()
if pioneer.get_nav_system().name != "GPS":
    raise SystemExit("Навигация не GPS")

lat, lon, alt = pioneer.get_global_position_gps()
pioneer.arm()
pioneer.takeoff()
pioneer.go_to_global_point(lat + 0.0001, lon, alt + 5)
while not pioneer.point_reached():
    pass
pioneer.land()
\end{verbatim}
\end{codebox}

\begin{warnbox}[Особенности глобальных точек в 0.11.0]
\begin{itemize}
    \item \verb|yaw=0| означает «курс не задавать»: библиотека выставляет флаг курса только при \verb|yaw != 0|. Нулевой курс (север) этой командой не задать; для него используйте \verb|set_yaw(0)| после прибытия.
    \item В исходнике метода \verb|go_to_global_point_relative| стоит пометка \texttt{CHECK}: смещения широты и долготы переводятся в целые числа иначе, чем абсолютные координаты, а описание единиц измерения высоты в исходнике и в README расходится (градусы или метры). Проверьте единицы на безопасной высоте малым смещением, прежде чем доверять команде.
    \item Пример с \verb|0.0001| градуса --- порядка 11~м по широте; большие смещения проверяйте на открытом месте.
\end{itemize}
\end{warnbox}

В симуляторе GPS и глобальные точки не поддерживаются: навигация там локальная (\verb|NavSystem.LPS|), поэтому проверка выше остановит программу.

\begin{taskbox}[Защита от неверной навигации]
Напишите функцию \verb|go_global(lat, lon, alt)|, которая сначала проверяет \verb|get_nav_system()| и \verb|get_nav_status_gps()|, и только при навигации GPS и статусе \texttt{OK} отдаёт команду, иначе печатает причину отказа и возвращает \verb|False|. В симуляторе вызовите её для любой точки.
\expected{в симуляторе функция печатает, что навигация не GPS, и возвращает \texttt{False}, не вызывая \texttt{go\_to\_global\_point}; на дроне с исправным GPS --- отправляет точку.}
\end{taskbox}

\subsection{Курс и возврат домой}

\begin{codebox}[set\_yaw и rtl]
\verb|set_yaw(yaw)| --- повернуться на месте на курс \verb|yaw| (градусы). Возвращает \verb|bool|. Действует только в воздухе, как полётная команда.\\
\verb|rtl()| --- вернуться домой (Return To Launch). Команда не блокирующая: возвращает \verb|True| сразу после отправки.
\end{codebox}

\verb|rtl()| не проверяет состояние и не ждёт окончания полёта. Дрон сам возвращается в точку взлёта; признак, что всё закончилось --- состояние \verb|ON_LAND|. Чтобы дождаться посадки, не вызывайте \verb|land()| (он ждёт события, которое придёт ещё не скоро), а опрашивайте состояние (имя состояния сравнивают через \verb|.name|, так работает и на дроне, и в симуляторе):

\begin{codebox}[Ждём возвращения]
\small
\begin{verbatim}
pioneer.rtl()
while pioneer.get_fly_state().name != "ON_LAND":
    time.sleep(0.2)
print("Дома")
\end{verbatim}
\end{codebox}

Состояние \verb|ON_LAND| SDK ставит по событию \verb|COPTER_LANDED|, так что цикл закончится, когда автопилот сообщит о посадке. Добавьте тайм-аут: если цикл ждёт дольше ожидаемого, вызовите \verb|land()| сами.

\begin{taskbox}[Домой по сигналу]
Взлетите, отлетите в точку $(1.5, 0, 1)$, затем вызовите \verb|rtl()| и дождитесь посадки циклом по \verb|get_fly_state()| с тайм-аутом 30~с. Напечатайте время возвращения.
\expected{дрон возвращается к точке взлёта и садится, программа печатает «Дома» и время; после выхода из цикла \texttt{get\_fly\_state()} равен \texttt{ON\_LAND}.}
\end{taskbox}

\subsection{События: подписка вместо опроса}

Опрос \verb|point_reached()| знает только про одну точку. Автопилот сообщает и о другом: батарея садится, сильный удар, моторы запущены. В SDK~2 на события можно подписаться --- как на \verb|callback| в Lua:

\begin{codebox}[subscribe и unsubscribe]
\verb|pioneer.subscribe(функция, Event.LOW_VOLTAGE1)| --- вызывать функцию при этом событии; без второго аргумента --- при любом событии. Функция получает событие одним аргументом; \verb|event.name| --- его имя.\\
\verb|pioneer.unsubscribe(функция, событие)| --- отписаться: передайте те же функцию и событие, что при подписке. Если событие в \verb|subscribe| не указывали, то и здесь не указывайте (по умолчанию --- \verb|Event.ALL|).
\end{codebox}

Одну функцию можно подписать на несколько событий, и на одно событие --- несколько функций. Список событий лежит в перечислении \verb|Event|; у каждого события есть числовой код автопилота:

\begin{center}
\small
\begin{tabularx}{\linewidth}{@{}>{\ttfamily}l r X@{}}
\toprule
\textrm{\textbf{Событие}} & \textbf{Код} & \textbf{Когда приходит} \\
\midrule
ALL & 255 & не событие: подписка «на всё» \\
ENGINES\_STARTED & 56 & моторы запущены \\
TAKEOFF\_COMPLETE & 51 & взлёт завершён \\
POINT\_DECELERATION & 43 & дрон начал тормозить перед точкой \\
POINT\_REACHED & 42 & точка достигнута \\
COPTER\_LANDED & 26 & дрон сел \\
LOW\_VOLTAGE1 & 31 & батарея садится, но вернуться домой хватит \\
LOW\_VOLTAGE2 & 32 & батарея почти пуста, автопилот начал аварийную посадку \\
LOW\_CHARGE & 115 & критически низкий заряд: запуск моторов невозможен \\
SHOCK & 65 & столкновение или сильный удар, возможна потеря управления \\
\bottomrule
\end{tabularx}
\end{center}

\begin{infobox}[Событие DISARM]
Справочник симулятора, который ссылается на версию 0.16, перечисляет ещё событие \verb|DISARM| («моторы выключены»). В исходниках 0.11.0 такого события нет. Если вашей версии оно неизвестно, подписка \verb|Event.DISARM| завершится ошибкой \verb|AttributeError|; проверяйте так: \verb|hasattr(Event, "DISARM")|.
\end{infobox}

\begin{codefile}[python]{python-onboard/pioneer_sdk2/sdk2_04_events.py}
\end{codefile}

\paragraph{Что происходит.} Это патруль с аварийным возвратом --- тот же автомат, что в Lua-миссии из главы~\ref{sec:missions-architecture}, только на Python. Функция \verb|on_event| печатает каждое событие, а при \verb|LOW_VOLTAGE1| поднимает флаг \verb|low_battery|. Решение принимает основной цикл: \verb|fly_to| замечает флаг, прерывает ожидание точки, и автомат переходит из \verb|PATROL| в \verb|RETURN|. Домой дрон летит уже без проверки флага (\verb|can_abort=False|): иначе низкий заряд прервал бы и возвращение.

\subsubsection*{Ждать событие, а не опрашивать}

\begin{codefile}[python]{python-onboard/pioneer_sdk2/09_events.py}
\end{codefile}

\paragraph{Что происходит.} Словарь \verb|EVENT_COLORS| связывает событие с цветом; подписка в цикле по ключам словаря избавляет от повторяющихся строк. \verb|on_event(event)| вызывается библиотекой в параллельном потоке. \verb|threading.Event| --- потокобезопасный флажок: обработчик делает \verb|set()|, основной поток ждёт \verb|wait(timeout)|. Порядок «\verb|clear()| $\to$ команда $\to$ \verb|wait()|» важен: если сбросить флажок после команды, можно пропустить быстро пришедшее событие. В \verb|finally| подписки снимаются через \verb|unsubscribe()|.

\begin{rulebox}[Обработчик события --- короткий]
SDK вызывает обработчик из своего служебного потока, который принимает данные от автопилота. Пока обработчик работает, следующие события ждут. Поэтому в обработчике только запоминают факт (флаг, счётчик), а летают --- в основном цикле. Блокирующую команду вроде \verb|land()| из обработчика не вызывайте: она ждёт события, которое этот же поток и должен доставить.
\end{rulebox}

\begin{taskbox}[Удар в патруле]
Добавьте в патруль обработку события \verb|SHOCK|: обработчик только поднимает флаг, а основной цикл, заметив его, включает красную подсветку и сажает дрон. Каждое событие печатайте по имени \verb|event.name|.
\expected{в обычном полёте в журнале \texttt{ENGINES\_STARTED}, \texttt{TAKEOFF\_COMPLETE}, \texttt{POINT\_REACHED}\ldots; при ударе --- \texttt{SHOCK}, красный свет и посадка из основного цикла.}
\end{taskbox}

\begin{taskbox}[Низкий заряд]
По событию \verb|LOW_VOLTAGE1| включайте мигание красным и вызывайте \verb|rtl()|. Посчитайте время между \verb|POINT_DECELERATION| и \verb|POINT_REACHED| для разных расстояний перелёта.
\expected{в журнале видна пара событий с интервалом, растущим с расстоянием; при низком заряде дрон мигает красным и возвращается домой.}
\end{taskbox}

\begin{taskbox}[Подсчёт точек]
Подпишите функцию на \verb|Event.POINT_REACHED| и считайте прибытия в счётчик; пролетите три точки подряд, дожидаясь прибытия в каждую по значению счётчика, а не через \verb|point_reached()|. В конце отпишитесь через \verb|unsubscribe| тем же способом, каким подписывались.
\expected{счётчик растёт на единицу за каждую точку и равен 3 к концу маршрута; после \texttt{unsubscribe} счётчик больше не меняется.}
\end{taskbox}

\subsection{Полёт со скоростью}

\begin{codefile}[python]{python-onboard/pioneer_sdk2/sdk2_05_manual_speed.py}
\end{codefile}

\paragraph{Что происходит.} В SDK~2 у \verb|set_manual_speed| есть аргумент \verb|interval|: одна команда держит скорость столько секунд, после чего дрон сам останавливается. Повторять команду в цикле каждые 0.05~с, как в главе~\ref{sec:python-manual-speed}, не нужно. В регуляторе, который пересчитывает скорость на каждом кадре, \verb|interval| делают коротким (0.3--0.5~с): если программа зависнет, дрон через полсекунды остановится сам. Скорость поворота \verb|yaw_rate| --- в рад/с и не больше 1.82; большее значение SDK отбросит без сообщения.

\subsubsection*{Повторять команду скорости}

\begin{codefile}[python]{python-onboard/pioneer_sdk2/08_manual_speed.py}
\end{codefile}

\paragraph{Что происходит.} Здесь команда повторяется каждые \verb|SEND_PERIOD| = 0.2~с, а \verb|COMMAND_INTERVAL| выбран чуть больше (0.3~с): между повторами нет «провалов», а если программа зависнет, дрон остановится через 0.3~с сам. В конце движения отправляется нулевая скорость --- остановка происходит сразу, а не после истечения интервала. Скорости можно комбинировать: \verb|vy| и \verb|yaw_rate| вместе дают полёт по дуге. Радиус дуги $R = v/\omega$.

\begin{codebox}[Команды скорости]
\verb|set_manual_speed(vx, vy, vz, yaw_rate, interval=1.0)| --- лететь со скоростью в осях системы координат (м/с); \verb|yaw_rate| в рад/с.\\
\verb|set_manual_speed_body_fixed(vx, vy, vz, yaw_rate, interval=1.0)| --- то же, но скорости заданы в осях самого дрона. В симуляторе \verb|vx| направлена вправо, \verb|vy| --- вперёд.
\end{codebox}

Команда скорости не ждёт и ничего не возвращает; скорость округляется до целых сантиметров в секунду, поэтому значения меньше 0.01~м/с превращаются в ноль. Остановить дрон раньше \verb|interval| можно командой с нулевыми скоростями.

\begin{taskbox}[Скорость без цикла]
Одной командой \verb|set_manual_speed(..., interval=2)| пролетите вперёд со скоростью 0.3~м/с, затем поверните на месте с \verb|yaw_rate=0.5| в течение 3~с. Сравните измеренные положение и курс \verb|get_local_yaw_lps()| до и после.
\expected{смещение около 0.6~м, после чего дрон сам зависает; курс меняется примерно на $86^\circ$ (0.5~рад/с $\times$ 3~с).}
\end{taskbox}

\begin{taskbox}[Плавный старт]
Сделайте «плавный старт»: скорость линейно растёт от 0 до 0.4~м/с за 2~с. Оцените радиус дуги по формуле $R = v/\omega$ и проверьте на практике. Реализуйте полёт по квадрату только командами скорости и сравните точность с примером 05.
\expected{дрон разгоняется без рывка; радиус дуги близок к $v/\omega$; квадрат из скоростей замыкается хуже, чем квадрат из точек.}
\end{taskbox}

\begin{taskbox}[Слишком быстрый поворот]
Подайте \verb|yaw_rate=2.0| на 2~с, затем \verb|yaw_rate=1.8| на 2~с, измеряя курс через \verb|get_local_yaw_lps()|. Объясните различие.
\expected{при \texttt{yaw\_rate=2.0} (больше 1.82) дрон не поворачивается, и ошибки нет --- SDK отбросил команду; при 1.8 курс меняется примерно на $206^\circ$ за 2~с (1.8~рад/с $\times$ 2~с).}
\end{taskbox}
